\section{Dependency Table}
\label{app:deps}

Table~\ref{tab:deps} lists, for each statement of this paper on the path to
Theorem~\ref{thm:main} whose proof cites another statement, the statements of
this paper and of Part~I that its proof uses and the parts of
Hypothesis~I.\ref{I-hyp:N} on which it depends,
directly or through its inputs. Statements whose proofs cite no other statement
have no row, and definitions are not listed as inputs. Theorem~\ref{thm:smooth}
and Lemma~\ref{lem:smooth-bounds} have a joint proof and share a row; consecutive
statements with the same inputs that depend on the same parts of
Hypothesis~I.\ref{I-hyp:N} also share a row. The rows are
ordered so that every input that has a row appears in an earlier row.
Statements of Part~I appear only as inputs, with their Part~I numbers; they
have no row, and a row is marked with the parts of Hypothesis~I.\ref{I-hyp:N}
that its Part~I inputs assume in their statements.

A statement that assumes some of \textup{(N1)}--\textup{(N3)} is marked with them,
and so is every statement whose inputs are marked. The objects of
Definition~I.\ref{I-def:eigendata} exist only under \textup{(N1)}--\textup{(N3)}, and
the seed of the approximate data of Definition~\ref{def:approx} is built from
them; so every statement about these objects, or about the approximate
solutions, depends on \textup{(N1)}--\textup{(N3)} through these definitions,
even when its proof cites none of the three. Theorem~\ref{thm:N-holds} proves
\textup{(N1)}--\textup{(N3)}, so it and Theorem~\ref{thm:main} depend on none of
them. A dash means that the statement depends on none of them.
Lemma~\ref{lem:ode} is quoted from \cite{Teschl} without proof, and only the proof of
Proposition~\ref{prop:runs} relies on machine computation.

\begingroup
\small
\captionsetup{type=table}
\caption{Dependencies of the statements on the path to Theorem~\ref{thm:main}, and the parts of Hypothesis~I.\ref{I-hyp:N} that they use.}
\label{tab:deps}
\setlength{\LTleft}{0pt}\setlength{\LTright}{\fill}\setlength{\LTpre}{0pt}
\begin{longtable}{p{0.25\textwidth}p{0.54\textwidth}p{0.12\textwidth}}
Statement & Inputs of the proof & Uses \\
\hline
\endfirsthead
\multicolumn{3}{l}{\emph{Table~\ref{tab:deps} (continued)}} \\
Statement & Inputs of the proof & Uses \\
\hline
\endhead
\raggedright Lem.~\ref{lem:flow}
  & \raggedright Lem.~\ref{lem:ode} & --- \\
\raggedright Lem.~\ref{lem:grad-bound}
  & \raggedright Lem.~I.\ref{I-lem:BS-basic} & --- \\
\raggedright Lem.~\ref{lem:comm-R2}
  & \raggedright Lem.~\ref{lem:embed} & --- \\
\raggedright Lem.~\ref{lem:iterates}
  & \raggedright Lem.~I.\ref{I-lem:BS-basic}, \ref{lem:flow} & --- \\
\raggedright Lem.~\ref{lem:iter-lip}
  & \raggedright Lem.~\ref{lem:flow}, \ref{lem:grad-bound}, \ref{lem:iterates} & --- \\
\raggedright Lem.~\ref{lem:iter-Hs}
  & \raggedright Lem.~I.\ref{I-lem:BS-basic}, \ref{lem:comm-R2}, \ref{lem:iterates}, \ref{lem:iter-lip} & --- \\
\raggedright Lem.~\ref{lem:iter-cauchy}
  & \raggedright Lem.~I.\ref{I-lem:BS-basic}, \ref{lem:iterates} & --- \\
\raggedright Prop.~\ref{prop:iter-limit}
  & \raggedright Lem.~I.\ref{I-lem:BS-basic}, \ref{lem:embed}, \ref{lem:iterates}, \ref{lem:iter-lip}, \ref{lem:iter-Hs}, \ref{lem:iter-cauchy} & --- \\
\raggedright Thm.~\ref{thm:smooth}; Lem.~\ref{lem:smooth-bounds}
  & \raggedright Prop.~\ref{prop:iter-limit}; Lem.~I.\ref{I-lem:BS-basic}, \ref{lem:flow} & --- \\
\raggedright Lem.~\ref{lem:Hk-calculus}
  & \raggedright Lem.~I.\ref{I-lem:BS-basic}, \ref{lem:embed} & --- \\
\raggedright Lem.~\ref{lem:duhamel}
  & \raggedright Lem.~I.\ref{I-lem:sector}, I.\ref{I-lem:semigroup-construct} & --- \\
\raggedright Lem.~\ref{lem:reduction}
  & \raggedright Thm.~\ref{thm:smooth}; Lem.~I.\ref{I-lem:ring-facts}, \ref{lem:flow} & --- \\
\raggedright Lem.~\ref{lem:bilinear}
  & \raggedright Lem.~I.\ref{I-lem:uniform}, \ref{lem:embed}, \ref{lem:Hk-calculus}, \ref{lem:reduction} & (N1)--(N3) \\
\raggedright Lem.~\ref{lem:energy}
  & \raggedright Prop.~I.\ref{I-prop:U}; Lem.~\ref{lem:embed}, \ref{lem:Hk-calculus}, \ref{lem:reduction} & (N1)--(N3) \\
\raggedright Lem.~\ref{lem:majorant}
  & \raggedright Prop.~I.\ref{I-prop:SU} & (N1)--(N3) \\
\raggedright Thm.~\ref{thm:quiet}
  & \raggedright Thm.~\ref{thm:smooth}; Prop.~I.\ref{I-prop:semigroup}; Lem.~\ref{lem:annuli}, \ref{lem:duhamel}, \ref{lem:reduction}, \ref{lem:bilinear}, \ref{lem:energy} & (N1)--(N3) \\
\raggedright Lem.~\ref{lem:running}
  & \raggedright Prop.~I.\ref{I-prop:semigroup} & (N1)--(N3) \\
\raggedright Lem.~\ref{lem:linear-part}
  & \raggedright Prop.~I.\ref{I-prop:semigroup}; Lem.~I.\ref{I-lem:uniform}, \ref{lem:running} & (N1)--(N3) \\
\raggedright Lem.~\ref{lem:window}
  & \raggedright Prop.~I.\ref{I-prop:semigroup}; Lem.~I.\ref{I-lem:uniform}, \ref{lem:annuli}, \ref{lem:duhamel}, \ref{lem:reduction}, \ref{lem:bilinear}, \ref{lem:energy}, \ref{lem:running}, \ref{lem:linear-part} & (N1)--(N3) \\
\raggedright Thm.~\ref{thm:active}
  & \raggedright Thm.~\ref{thm:smooth}; Prop.~I.\ref{I-prop:semigroup}; Lem.~I.\ref{I-lem:uniform}, \ref{lem:annuli}, \ref{lem:duhamel}, \ref{lem:reduction}, \ref{lem:bilinear}, \ref{lem:energy}, \ref{lem:running}, \ref{lem:linear-part}, \ref{lem:window} & (N1)--(N3) \\
\raggedright Cor.~\ref{cor:active-moment}
  & \raggedright Thm.~\ref{thm:active}; Lem.~I.\ref{I-lem:uniform} & (N1)--(N3) \\
\raggedright Lem.~\ref{lem:seed}
  & \raggedright Thm.~I.\ref{I-thm:evans-simple}; Lem.~I.\ref{I-lem:sector}, \ref{lem:flow} & (N1)--(N3) \\
\raggedright Prop.~\ref{prop:integrability}
  & \raggedright Lem.~I.\ref{I-lem:ring-facts} & --- \\
\raggedright Lem.~\ref{lem:adm}
  & \raggedright Thm.~\ref{thm:smooth}; Prop.~\ref{prop:integrability}; Lem.~I.\ref{I-lem:ring-facts}, \ref{lem:seed} & (N1)--(N3) \\
\raggedright Lem.~\ref{lem:regions-basic}
  & \raggedright Thm.~\ref{thm:smooth}; Lem.~I.\ref{I-lem:ring-facts}, \ref{lem:annuli}, \ref{lem:seed}, \ref{lem:adm} & (N1)--(N3) \\
\raggedright Lem.~\ref{lem:inner-circ}
  & \raggedright Thm.~\ref{thm:smooth}; Lem.~\ref{lem:regions-basic} & (N1)--(N3) \\
\raggedright Lem.~\ref{lem:speed}
  & \raggedright Lem.~I.\ref{I-lem:BS-basic}, I.\ref{I-lem:ring-facts}, \ref{lem:seed}, \ref{lem:inner-circ} & (N1)--(N3) \\
\raggedright Lem.~\ref{lem:clock}
  & \raggedright Prop.~I.\ref{I-prop:SU}; Lem.~\ref{lem:inner-circ} & (N1)--(N3) \\
\raggedright Lem.~\ref{lem:transfer}
  & \raggedright Lem.~\ref{lem:hollow}, \ref{lem:ext-multipole}, \ref{lem:regions-basic}, \ref{lem:speed} & (N1)--(N3) \\
\raggedright Lem.~\ref{lem:inner-field}
  & \raggedright Lem.~\ref{lem:ext-multipole}, \ref{lem:regions-basic}, \ref{lem:inner-circ}, \ref{lem:transfer} & (N1)--(N3) \\
\raggedright Prop.~\ref{prop:confine}
  & \raggedright Thm.~\ref{thm:smooth}; Lem.~\ref{lem:hollow}, \ref{lem:regions-basic}, \ref{lem:speed} & (N1)--(N3) \\
\raggedright Lem.~\ref{lem:field-bound}
  & \raggedright Lem.~\ref{lem:regions-basic}, \ref{lem:speed}, \ref{lem:transfer}, \ref{lem:inner-field} & (N1)--(N3) \\
\raggedright Lem.~\ref{lem:coupling}
  & \raggedright Thm.~\ref{thm:quiet}; Prop.~I.\ref{I-prop:semigroup}; Lem.~I.\ref{I-lem:BS-basic}, I.\ref{I-lem:semigroup-construct}, I.\ref{I-lem:uniform}, \ref{lem:reduction}, \ref{lem:regions-basic} & (N1)--(N3) \\
\raggedright Lem.~\ref{lem:detuning}
  & \raggedright Prop.~I.\ref{I-prop:U}; Lem.~\ref{lem:inner-circ} & (N1)--(N3) \\
\raggedright Lem.~\ref{lem:transfer-step}
  & \raggedright Lem.~\ref{lem:ext-multipole}, \ref{lem:regions-basic}, \ref{lem:inner-circ}, \ref{lem:transfer} & (N1)--(N3) \\
\raggedright Lem.~\ref{lem:O}
  & \raggedright Thm.~\ref{thm:quiet}; Lem.~\ref{lem:reduction}, \ref{lem:majorant}, \ref{lem:detuning}, \ref{lem:transfer-step} & (N1)--(N3) \\
\raggedright Prop.~\ref{prop:relay}
  & \raggedright Thm.~\ref{thm:quiet}, \ref{thm:active}; Lem.~\ref{lem:majorant}, \ref{lem:clock}, \ref{lem:transfer}, \ref{lem:detuning}, \ref{lem:transfer-step}, \ref{lem:O} & (N1)--(N3) \\
\raggedright Lem.~\ref{lem:Oprime}
  & \raggedright Thm.~\ref{thm:quiet}, \ref{thm:active}; Prop.~\ref{prop:relay}; Lem.~\ref{lem:detuning}, \ref{lem:transfer-step}, \ref{lem:O} & (N1)--(N3) \\
\raggedright Lem.~\ref{lem:drift}
  & \raggedright Thm.~\ref{thm:active}; Prop.~\ref{prop:confine}; Lem.~I.\ref{I-lem:uniform}, \ref{lem:running}, \ref{lem:inner-circ}, \ref{lem:coupling}, \ref{lem:detuning} & (N1)--(N3) \\
\raggedright Lem.~\ref{lem:handoff}
  & \raggedright Thm.~\ref{thm:quiet}; Prop.~I.\ref{I-prop:semigroup}; Lem.~I.\ref{I-lem:uniform}, \ref{lem:majorant}, \ref{lem:inner-circ}, \ref{lem:coupling} & (N1)--(N3) \\
\raggedright Lem.~\ref{lem:kicktrue}
  & \raggedright Thm.~\ref{thm:quiet}, \ref{thm:active}; Prop.~\ref{prop:confine}, \ref{prop:kick}; Lem.~I.\ref{I-lem:ring-facts}, I.\ref{I-lem:uniform}, \ref{lem:annuli}, \ref{lem:majorant}, \ref{lem:inner-circ}, \ref{lem:clock}, \ref{lem:coupling}, \ref{lem:detuning}, \ref{lem:Oprime}, \ref{lem:drift}; Cor.~\ref{cor:active-moment} & (N1)--(N3) \\
\raggedright Lem.~\ref{lem:time}
  & \raggedright Lem.~\ref{lem:detuning} & (N1)--(N3) \\
\raggedright Lem.~\ref{lem:quiet-time}
  & \raggedright Lem.~\ref{lem:time} & (N1)--(N3) \\
\raggedright Lem.~\ref{lem:theta-quiet}
  & \raggedright Lem.~\ref{lem:O}, \ref{lem:quiet-time} & (N1)--(N3) \\
\raggedright Lem.~\ref{lem:nonempty}
  & \raggedright Thm.~\ref{thm:quiet}, \ref{thm:active}; Prop.~I.\ref{I-prop:U}; Lem.~\ref{lem:inner-circ}, \ref{lem:Oprime} & (N1)--(N3) \\
\raggedright Lem.~\ref{lem:geometry}
  & \raggedright Prop.~\ref{prop:confine}; Lem.~I.\ref{I-lem:uniform}, \ref{lem:reduction}, \ref{lem:clock}, \ref{lem:field-bound}, \ref{lem:detuning} & (N1)--(N3) \\
\raggedright Prop.~\ref{prop:bootstrap}
  & \raggedright Thm.~\ref{thm:smooth}, \ref{thm:quiet}, \ref{thm:active}; Prop.~\ref{prop:confine}, \ref{prop:relay}; Lem.~I.\ref{I-lem:uniform}, \ref{lem:reduction}, \ref{lem:seed}, \ref{lem:regions-basic}, \ref{lem:speed}, \ref{lem:clock}, \ref{lem:inner-field}, \ref{lem:coupling}, \ref{lem:detuning}, \ref{lem:O}, \ref{lem:Oprime}, \ref{lem:drift}, \ref{lem:handoff}, \ref{lem:kicktrue}, \ref{lem:time}, \ref{lem:quiet-time}, \ref{lem:theta-quiet}, \ref{lem:nonempty}, \ref{lem:geometry}; Cor.~\ref{cor:active-moment} & (N1)--(N3) \\
\raggedright Prop.~\ref{prop:enclosure}; Lem.~\ref{lem:apriori}
  & \raggedright Lem.~\ref{lem:end-regions} & --- \\
\raggedright Lem.~\ref{lem:Nprime-N}
  & \raggedright Lem.~I.\ref{I-lem:evans-basic}, \ref{lem:apriori} & (N1) \\
\raggedright Prop.~\ref{prop:runs}
  & \raggedright Prop.~\ref{prop:enclosure}, \ref{prop:winding}; Lem.~I.\ref{I-lem:evans-basic}; computation under (TB1)--(TB3) & --- \\
\raggedright Prop.~\ref{prop:N3}
  & \raggedright Lem.~\ref{lem:apriori} & --- \\
\raggedright Prop.~\ref{prop:reduction}
  & \raggedright Thm.~\ref{thm:active}; Prop.~\ref{prop:integrability}, \ref{prop:bootstrap}; Lem.~I.\ref{I-lem:radial-stationary}, I.\ref{I-lem:uniform}, \ref{lem:adm}, \ref{lem:inner-circ}, \ref{lem:time}, \ref{lem:nonempty}; Cor.~I.\ref{I-cor:defect} & (N1)--(N3) \\
\raggedright Thm.~\ref{thm:N-holds}
  & \raggedright Prop.~\ref{prop:runs}, \ref{prop:N3}; Lem.~\ref{lem:Nprime-N} & --- \\
\raggedright Thm.~\ref{thm:main}
  & \raggedright Thm.~\ref{thm:N-holds}; Prop.~\ref{prop:reduction} & --- \\
\end{longtable}
\endgroup
